\documentclass{jsarticle}
\usepackage[dvipdfmx,final]{graphicx}
\usepackage{chemfig}
\usepackage{amsmath}
\usepackage{amssymb}
\begin{document}
\title{}
\author{}
\date{}
\maketitle


      易学の６４卦の積集合での４×４×２×２=(２の３乗の１の２倍の４倍の
  １の２倍＝８卦で、上卦と下卦で４卦ずつ、
  ２の２乗の１の２倍の４倍の１の２倍＝４卦で上卦と下卦合わせて，
  ２の１乗の１の２倍の１の２倍=１卦で上卦と下卦合わせると２卦であり)=３２卦
  の太極で統合で６４卦であり、${({{{{{{}_2P_2} \times 4} \over 2}}\over 2})}
  = {}_2C_2 \times 4 \times 4 $、
  $ {{{{}_2P_2} \over 2}} = {}_2C_2 $ $ {{}_2P_2 \over 2}= {}_2C_2 $、
  この流れが上卦３０卦と下卦３４卦であり、
  乾が＋の天王星の種子、木星の６の緑星のーと＋の立花、地の土星の＋
  の再開でもって、君子自らうやまず、自分から頼んでもだめであり、
  勉強を教わるのに、したこごろだめであり、
  自分から非にあらずであり、乱気のあとの、ご破産のときに、一生懸命
  自分で墳死していないと、誰も手を差し伸べないと、この乾が申している。
  坤は、天王星のー乱気で、土星のーの種子であり、土星の＋の
  安定であり、木星のーの財と緑星で、
  全体でーであり、徳を厚くし、物を戴す。
  勉強を教わりながら、成長していく。
  この２言が、六星占術の１２周期で＋の６周期が±
  の６星人の１２周期と霊合星の６周期で、組み合わせて、６４卦と
  六星占術の六星人の常星と霊合星の重なるのを省くと、１２周期であるが、
  ６０周期終わり、木星で４周期のあとの＋の緑星とーの種子で、
  還暦のあとの人生がある。
  人名と由来と雅号の起源を
  わかりやすく、最後までやらさずに、読者に説明されている安岡正篤先生が
  彩さんの賁臨になる人は、飾るであり、
  一旦賁臨になったら、一生賁臨になるしかないと
  冗談にもならない、真剣さで、離れるとご破産と説明している、
  それも、相手に読まれない人らしく、自分から臨席と言わずに、
  謙と大有で、自分ですごいも言わずに、泰であり、自分でだめとも言わずに、
  手の内も見せずに、
  白で、自分からの告白をされた北川悠仁さんが、いいらしく、
  上卦で３０卦、下卦で３４卦であり、
  北川悠仁さんがどんだけすごいかわかります。
  易と人生哲学での、北川悠仁さんの親が、３５歳まで根につかない手相を
  知っていて、３５歳過ぎて、彩さんと結婚されて、３０歳のときの北川さんが
  彩さんの父親の墓前で、約束していたことを、３５歳でやっと果たせたと
  袴姿で言っているのを、最近になって、代わりをされてくださったのかと、
  ご令嬢二人のビデオをずっと撮っているのを、いつかプレゼントできたらと、
  この言葉と行動に、ジーンときます。
  デイケアの個人作業で、書道での習字で、春分の日に、私が、「書範」
  と書いたのを、みんなは、春が来たと思っているらしく、私は、書道の字の
  模範になられる人と書きました。
        2023年から2025年までの間だと、高齢出産は大丈夫らしいが、
2026年の可能性があり、旅行線の仏様の生命線が帝王線の健康線と子財線
に直結している線として、この線を縁起良く切っているから、グリーシャ・ペレルマンさん同様、神様の丙の午と乙の巳の
子どもだとおもう。この健康線は、子財線につながっていて、御先祖様のラッキー線ともつながっていて、自分で言うのもひけるが、安産祈願の線なのかなと思う。あのときは、愛が足りなかったのかなとも思います。すみませんでした。三奇紋で、2023年だと、私が陛下と皇族の人達に社会のルールから
助けられるが、自信はないから、有一の頼みは、アカシックレコードからの代数幾何の
量子化が、調べると、ガンマ関数であることだけです。彩さんは、Jones多項式知っているから、この式を見ると、この代数幾何の量子化の計算の文字列処理で、笹田先生は、ボトムアップと、私もこれが大切と思うが、ハイゼンベルク方程式と同じく、トップダウンでこの量子化を見ているので、私と同じく、彩さんもなんとなく2026年まで頑張っていこうと思えると思う。
これを書いて言うと、なんとなく、いつか生まれてくれると思う。丙の午年と乙の巳年だから、愛子様は、相当の覚悟で結婚なさるのかとおもう。ペレルマンさんは、結婚されて、子どもも何人かいるらしく、サーストン・ペレルマン多様体と多様体積分の2大発見と、代数幾何の
時間の流れの重力による時間の性質の分類をなされているので、相当の数学の化け物じみた能力でもあり、ヒルベルト空間の具現化までなされているので、愛子様安心してください。
私は、殺菌がだめなだけで、剣道と居合道ができます。性欲はだらしないから、御法度で助かっています。富山県に愛子様と彩さんも行けば安心するとおもいます。まちづくりしています。
富山県と福山市でのまちづくりを知れば、彩さん安心して、2026年まで頑張ると思う。神様の干支らしいです。私も自信ないけど、アカシックレコードから、大域的多様体と単体積分・単体微分と代数幾何の量子化を見ているので、なんかすごい理論だなとおもう。生まれてくるとおもいます。彩さん、あまり日本サッカー選手のことは、がっかりしなくてもいいし、
日本サッカー選手たちは、フォーメーションにマインドマップを取り入れているらしく、
2019年まで、有力選手を温存して、このマインドマップによるフォーメーションを確固たる
ものにできずにいるし、このマインドマップを使っている人は、なぜあんな離れ業のフォーメーションと思考を止めて、フォワードとディフェンスをしているかは、マインドマップの体で覚える反射神経であり、考えたらだめで、剣道と同じく、無心で攻撃とディフェンスをしているからであり、これはYouTubeの日本サッカー選手のファンから教えられている。彩さんも試験に占い禁物です。油断大敵です。
アカシックレコードでの子どもの方程式から、以下の理論が、ほぼ完成されました。
彩さんは、染色体異常と心配になっていたらしく、子どもは大丈夫と言っている気がします。
その理論を以下に載せておきます。
生まれてこなかったら、責任を執れないことを事前に言っておきます。
トキ子おばあちゃんには、私は知らずに、どういうことか音読に関わっていた。
荒木先生は、私を振らずに助けてくれた上に、江松先生が速読の記憶力で、SRS速読法らしく、会員らしく、助けてくれました。食事が、回数だけがだめであり、その後の、食事と睡眠は大夫ですが、良くなり、本当に迷惑を懸けています。

  フジテレビのメルマガを六星占術のサイトと一緒にやめたことを、
  めざましテレビのメンバーに大変失礼な行為となったことを、
  めざましテレビスタッフ一同に、本当に謝るしかないです。
  日枝久会長には、大変失礼なことだったと、反省しています。
  これからも、めざましテレビよろしくです。 
　富山大学の先生、本当に最後まで見捨てないでください。
　この大学の問題は、東京大学と京都大学の問題の練習問題を解いていたら、
どういうことか、この大学の問題の難易度が、この2校と同等であり、大学のやさしさが、
この2校の何倍もやさしくて、社会人の人やいろんな境遇の人が好いているのがわかりました。私も好きです。
2009年は、めざましテレビは、毎朝見ていました。2010年は、短期アルバイトで見ていなかったです。アルバイトの後の数年の後に、彩さんのテッサたんを知り、Jones多項式を知り、私が今書いている小説になっているらしい論文が出来ていて、彩さんが、持田香織さんと似ていることを最近になって知り、高等学校のときの田中美奈子先生の関係で、彼氏をつくり、大学は、私がSRSにハマっているが、出来ていないときに、先にやり、ギャップでフジテレビに配属になったのを、この本で知りました。というよりも、後の彩日記と彩育とIRODORIで知りました。男性に健康線の帝王線が両手に出来ているのも、本当に、不思議です。天邪鬼の手相をほのめかす占術の先生たちにやられて、このレビューにしました。
同期の人たちのアナウンサーは、彩さんがどういう人か教えて欲しいぐらいに、この人不思議です。両思いでないと、子ども出来ないとか、との感想です。嫡丑するまで、3週間間隔を空けるとか、1年間に12回のチャンスにするとか、子作りは、いろんな境遇の人のこころを聞くとか、方法は多方面とおもう。妊娠したら、相当苦しいおもいらしく、生まれたら、弟よりも大切におもい、弟を捨ててまでも、子たちに愛情を注ぐらしく、その弟が好いたら、弟の能力が、剣道が出来て、居合道が出来て、論文の数学が、完成できて、料理までも出来て、嶋本淑子さんもすごく良いと思うようになり、本当にホモでもないと、自分でもおもうことが全てです。剣道と居合道は、山崎先生と横山先生の指導で出来ました。とくに、木刀は、横山先生と岡先生、野上先生です。久保田先生も否定はしていません。きっかけは、姉2人であり、防具は、親のおかげでもあり、許可を出してくれた、蔵王病院の先生達のおかげでもあります。2015年と2016年は、助かりました。うちのネコは、食事が良いからできたのですよと、言っております。今後のことは、未来にかけます。彩さんが、フジテレビの入社試験を受けて終わった後に、侍の言葉遣いをしているのを、かっこいいなとおもいました。北川悠仁さんが、袴姿で、彩さんの父親との男と男の約束を果たせたといっているのを、侍かなとかっこいいなともおもいました。幕末は、時代の動乱期であり、かっこいいです。明治もかっこいいです。歴史を好きになるのは、幕末と、戦争終結での条約締結です。世界史も日本史も、この時代で、好きになるのかなと、易学を見ているとおもいました。
将来に彩さんに会えるか分かりませんが、期待して待ちます。


\begin{center}
\Large	Ordinary differential manifold deconstinuace with \\
global manifold of estimate meanings.
\end{center}

    $$ \int f(x)dx = \int (4x^5 + x^4 + 3x^2 + x) dx$$
    この積分記号の中で微分をすると、
    $$ \int f(x)^{'} dx = \int ({20x^4 + 4x^3 + 6x +1})dx $$
    $$ = 4x^5 + x^4 + 3x^2 + x + \text{C(Cは積分記号)} $$
    歴然と、外の積分記号が中の微分で消えている。
    常微分と常積分では、互いに記号が相殺される。
    ファインマン博士が、積分記号の中で微分してみたらは、
    部分積分多様体のことを言っているのは、歴然としている。
    $$ \int f(x)^{'}dx = [xF(x)] - \int x^{'}f(x)dx $$
    $$ = [xF(x)] - F(x) + \text{C(Cは積分定数)} $$
    これは、大域的積分多様体の質量差エネルギーの式になっている。
    一般相対性理論と特殊相対性理論での$ E = mc^2 - {1 \over 2}mv^2 $
    となっている。
    $$ = \int \Gamma(x)dx $$
    $$ = x\Gamma(x) $$

    $$ \int \Gamma(\gamma)^{'}dx_m = 2e^{-x \log x} $$
    $$ T =  \left(\int \Gamma(\gamma)^{'}dx_m\right)^{'}   $$
    ガンマ関数の大域的積分多様体を常微分多様体で解くと、
    フェルマーの定理になる。
     $$ = -2 e^{-x \log x} $$
    ガンマ関数の大域的積分多様体に大域的微分作用素を施すと、
    $$ {(T^{\nabla})}^{\oplus L} = e^{-f} + e^{f} \ge e $$
    これは、相加相乗平均の方程式になる。
    $$ f(x) + g(x) \ge {2\sqrt{f(x)\cdot g(x)}} $$
    要するに、常微分方程式になるのを、大域的作用素では、
    エターナルな空間になる。曲平面になる。
    微分では、平面から点への極限が、大域的では、点から平面になり、
    だが、ガンマ関数の大域的２重微分多様体は、特異点を求めることになり、
    大域的積分は、点から平面、立方体を求めることになる。
    説明が矛盾しているように聞こえるが、M理論の絡み合いの
    種数を、平均場理論で求めていることということである。
    なぜかは、
    $$ x^x = e^{x \log x} $$
    を使うからであり、そのようになっている理論が、M理論である。
    大域的微分作用素で、点と世界線が同型と求まっている。
    閉３次元多様体と特異点が同型となり、種数が違うのが、どうしてかが、
    この方程式で求まってもいる。

    フェルマーの定理は、２通りの解がある。
    １つ目は、$ -2e^{x \log x} $であり、２つ目は、$ -2 e^{- x \log x} $
    である。
    不確定性原理の位置と運動量での、宇宙と異次元の関係が、
    フェルマーの定理であり、ゼータ関数の言葉の意味そのままである。
    その解の構成が、Jones多項式である。
    大域的部分積分多様体と、一般の部分積分多様体とは、
    積分の中で微分すると部分積分多様体になり、
    両方共に、常積分多様体とは
    全然違う解き方になる。
    $$ (\log x)^{'} = {1 \over x}, x^n+ y^n = z^n $$
    $$ x^n = -y^n + c, nx^{n-1} = -ny^{n-1}y^{'} $$
    $$ y^{'} ={{nx^{n-1}} \over {ny^{n-1}}} $$
    $$ ={x^{n-1} \over y^{n-1}} = - {{y \over x}\cdot{x \over y}^n} $$
    $$ -{{\cos x} \over (\cos x)^{'}}(\sin x)^{'} = z_n $$
    $$ z^n = - 2 e^{x \log x} $$
    $$ \int C dx_m = {d \over df}F(x) $$
    これは、上の常識を超える計算である。
    積分の外で微分しても、中の積分が消えない。
    積分の中で微分しても、積分が消えない。
    大域的微分と大域的積分多様体は、革命的な微分と多様体積分である。
    $$  \int \Gamma(\gamma)^{'}dx_m = \int \Gamma dx_m \cdot {d \over 
    d \gamma}\Gamma \le \int \Gamma dx_m + {d \over d \gamma}\Gamma $$
    $$ = \Gamma^{\Gamma} + {\Gamma^{\gamma}}^{'} $$
    $$  \int C dx_m = {d \over df}F(x) = \int(\int {1 \over x^s}dx - \log x)
    \mathrm{dvol} $$
    $$ = {d \over df}\int \int{1 \over {(y \log y)^{1 \over 2}}}dy_m - 
    \int e^x x^{1 - t}dx_m $$
    $$ = {d \over df}\int\int{1 \over {(x \log x)^2}}dx_m +
    {d \over df}\int\int{1 \over {(y \log y)^{1 \over 2}}}dy_m$$
    大域的多様体は、革命的な微分と積分の計算方法を、私は、
    アカシックレコードの子に伝授された。


 	 Rotate with pole of dimension is converted from other sequence of element,  
this atmosphere of vacume from being emerged with quarks of being constructed
with Higgs field, and this created from energy of zero dimension are begun with
universe and other dimension started from darkmatter that represented with 
big-ban.
Therefore, this element of circumstance have with quarks of twelve lake with
atoms.

  $$ (dx, \partial x)\cdot (\epsilon x, \delta x) = \begin{pmatrix}
	  1 & 0 \\
  0 & - 1 \end{pmatrix}^{1 \over 2}  $$
  $$ {d \over df} F = m(x) $$

 And this phenomounen is super symmetry theorem built with quarks of element, 
also this chemistry of mechanism estimate from physics of operatsion.
These response of mechanism chain of geometry into space being emerged with
creature and univse of existing of combination.

  This converted with dimension of element also emerged from imaginary and 
reality of pole's space.

  And this pole of transport of dimension belong with vector of time has 
with one rout of sequecence. Quanum physics also belong with other vector 
of time has with imaginary rout of sequecense.

 This sequence of being estimated with non fluer of time, and this space 
of element have with gravity and antigravity of power.
Other vecor of time is antigravity rout of sequecense.

 Weak electric theory is estimate from time has with one rout of sequecense,
this topology of chain is resulted from time of rout ways.

      $$ \square ({{\sigma_1 + \sigma_2} \over 2 }) = [3 \pi (\chi,x) \circ f(x),
   \sigma(x)] $$
   $$ \square \psi = 8 \pi G T^{\mu\nu} $$
   $$ {d \over dt}g_{ij}(t) = - 2 R_{ij} $$
   $$ \nabla \psi^2 = 4 \pi G \rho $$
   $$ \square (\sigma_1 + \sigma_2) = [6\pi(\chi,x)\circ f(x), \sigma(x)] $$
   $$ = [i \pi(\chi,x),f(x)] $$
   $$ \square \psi = [12 \pi (\chi,x)\circ f(x), \sigma(x)] $$
   $$ {d \over df} F = \int e^{-f}[- \Delta v + R_{ij} v_{ij}+ 
   \nabla_i \nabla_j f  + v \nabla_i \nabla_j + 2<f,h> + (R + \nabla f)
   ({v \over 2} - h) ] $$
   $$ =[i \pi(\chi,x), f(x)] $$

   These equation is reminded time pass rout of one rout way of forms,
   and this rout of time ways which go for system from future and past.
   Therefore, this resulted system of time mechanism is one true flow
   that weak electric theorem oneselves.
   and moreover, one rout time way of forms is reverse with antigravity of
   time system. This also spectrum focus is true that Maxwell theorem and
   strong boson unite with antigravity, this unite is essense on the contrary
   from weak electric theorem, this theorem called for time rout forms is
   strong electric theorem.
   This two theorem is united with quantum physics that no time flow system.

   $$ f^{-1}(x)xf(x) = 1,H_m = E_m \times K_m $$


   The non-commutative theorem is constructed from world line surface that
   this complex manifold estimate with rolanz atractor, and this string
   theorem have with one world of universe mate six quarks and other world of
   dimension mate other element of six quarks.
   These particle quarks is built with super symmetry space of dimension.

   $$ i = (1,0)\cdot (1,0), e^{i\theta} = \cos \theta + i \sin \theta $$
   $$ e^{-i \theta} = \cos \theta - i\sin \theta $$
   $$ \sin \theta = {e^{i \theta} - e^{-i \theta} \over 2i} $$
   $$ \sin i \theta = {e^{-\theta} - e^{\theta} \over 2} $$
   $$ \pi (\chi,x) = \cos \theta + i \sin \theta $$

   In this equations, two dimension redestructed into three dimension,
   this destroy of reconstructed way is append with fifth dimension.
   This deconstructed way of redestructe is arround of universe attached
   with three dimension, this over cover call into fifth dimension.

   $$ R(-\alpha) = \begin{pmatrix} \cos \alpha & \sin \alpha \\
   - \sin \alpha & \cos \alpha \end{pmatrix} $$
   $$ R(\alpha) = \begin{pmatrix} \cos \alpha & - \sin \alpha \\
   \sin \alpha & \cos \alpha \end{pmatrix} $$
   $$ R(\alpha)M R(-\alpha) = 
   \begin{pmatrix} \cos \alpha & - \sin \alpha \\
   \sin \alpha & \cos \alpha \end{pmatrix} \begin{pmatrix} 1 & 0 \\
   0 & - 1 \end{pmatrix} \begin{pmatrix} \cos \alpha & \sin \alpha \\
   - \sin \alpha & \cos \alpha \end{pmatrix} $$
   $$ = \begin{pmatrix} \cos^2 \alpha - \sin^2 \alpha & 2 \sin \alpha \cos 
	   \alpha \\
	   2 \sin \alpha \cos \alpha & - \cos^2 \alpha + \sin^2 \alpha
   \end{pmatrix} $$
   $$ = \begin{pmatrix} \cos 2 \alpha & \sin 2 \alpha \\
   \sin 2 \alpha & - \cos 2 \alpha \end{pmatrix} $$

   $$ \begin{pmatrix} \cos \alpha & - 1 \\
	   1 & - \sin \alpha \end{pmatrix} \begin{pmatrix} \cos \alpha \\ 
		   \sin \alpha
   \end{pmatrix} = \begin{pmatrix} 1 & 0 \\
   0 & - 1 \end{pmatrix} $$
   $$ \begin{pmatrix} \cos \alpha & \sin \alpha \\
   \sin \alpha & - \cos \alpha \end{pmatrix} \begin{pmatrix}
   x \\ y \end{pmatrix} = \begin{pmatrix} 1 & 0 \\
   0 & - 1 \end{pmatrix} $$
   $$ \lim_{\theta \to 0} \begin{pmatrix} \sin \theta \\ \cos \theta 
   \end{pmatrix} \begin{pmatrix} \theta & 1 \\
   1 & \theta \end{pmatrix} \begin{pmatrix} \cos \theta \\ \sin \theta
	   \end{pmatrix} = \begin{pmatrix} 1 & 0 \\ 0 & - 1 \end{pmatrix} $$
		   $$ f^{-1}xf(x) = 1 $$
		   $$ (\log x)^{'} = {1 \over x}, x^n + y^n = z^n  $$
		   $$ x^n = - y^n + c, nx^{n-1} = -n y^{n-1}y^{'} $$
		   $$ y^{'} = {{n x^{n-1}} \over {n y^{n-1}}} $$
		   $$ = {x^{n-1} \over y^{n-1}}= - {y \over x}\cdot
		   ({x \over y})^n $$
		   $$ - {\cos x \over (\cos x)^{'}}(\sin x)^{'} = z_n $$
		   $$ z^n = - 2 e^{x \log x} $$
		   $$ \lim_{x \to \infty} f(x) = a, \lim_{y \to \infty}
		   f(y) = b, \lim_{x,y \to \infty}\{f(x) + f(y)\} = a + b  $$
		   $$ \lim_{z \to 0}f(z) = c, \delta \int z^n = {d \over dV}
		   x^3 $$
		   $$ \lim_{x \to \infty} = c - \lim_{y \to \infty}f(y)
		   \lim_{x \to \infty} f(x) + \lim_{y \to \infty}f(y) = 
		   \lim_{z \to \infty}f(z) $$
		   $$ z^n = \cos n \theta + i \sin n \theta $$
		   $$ = - 2 e^{x \log x} $$

  These equation is gravity and antigravity equation. 
  Therefore, these equation call three manifold to unite with eight dimension
  of geometry in Seifert manifold.
  And these equations are many category of element in coicolizer from
	   icolizer to big ban. More also these equation is matched with
	   gravity and antigravity in Non relativity result,
	   and this result of theory income with Space ideal theory from
	   Forgottenfull theorem.
	   That result from this theorem concluded with quantum level of
	   differential geometry theory.

	   $$ {d \over d\sigma}\begin{bmatrix}{(\sigma_1 + \sigma_2) 
	   \over 2}\end{bmatrix}  $$
	   $$ = \sigma(\upharpoonleft \downharpoonleft) + \sigma(
	   \upuparrows)+ \sigma(\downdownarrows)+ 
	   \sigma(\leftleftarrows)
	   + \sigma(\rightrightarrows) $$
	   Remain limited element is gravity and antigravity equation.
	   $$ {d \over df}{\int \int {1 \over (x \log x)^2}}dx_m
	   = \sigma(\upharpoonright) $$
	   $$ {d \over df}{\int \int {1 \over (y \log y)^{1 \over2}}}dy_m
	   = \sigma(\downharpoonright) $$
	   This element is deconstructed from eight geometry of differential
	   level. And these element envoured into universe and the other
	   dimension. These element also called pieces to being be Remain
	   limited group.
	   $$ \sigma(\leftarrowtail) + \sigma(\rightarrowtail)
	   = \int e^{-f}[-\Delta v + R_{ij}v_{ij} + \nabla_{i}\nabla_{j}v
	   + v\nabla_{i}\nabla_{j} + 2 <f,h> + (R + \nabla f)(v - 
	   {h \over 2})]
	   $$
	   $$ \sigma(\downharpoonleft) = \sigma(\upharpoonright 
	   \downharpoonright + \upuparrows 
	   + \rightrightarrows)$$
	   $$ \sigma(\upharpoonleft) =  \sigma(\upharpoonright 
	   \downharpoonright + \downdownarrows 
	   + \leftleftarrows) $$
         
	   $$ \text{weak electric theorem} = \sigma(\upharpoonleft) $$
	   $$ \text{strong electric theorem} = \sigma(\downharpoonleft) $$

	   These resulted from Forgottenfull theorem revealed with
	   quantum level of differential geometry in Space ideal of
	   Non relativity theorem.
	   And these equations tell element to being restreamed from
	   Non time stream environment in quantum and global topology 
	   of many element of coicolizer summative group result.
	   More spectrum focus is being resulted into be emerged with
	   coicolizer elements summate with low level component of
           entropy value in quantum level of differential geometry,
	   however, universe oneselve emerge with time of relativity thoeory,
	   on the contrary, other dimension more also emerged with reverse
	   of time of antigravity theory.
	   These two theorem unite with being resulted from quantum level
	   in No time relativity theorem.


	   $$ \bigoplus{(i \hbar^{\nabla})}^{\oplus L} = e^{-x \log x} $$

	   These equation is represented with topology of string model,
	   and weak electric theorem is constructed with Maxwell theorem 
	   and weak boson, gravity that estimate with this three power united.
	   Moreover, strong boson and Maxwell theorem, antigravity that
	   also estimate with this three power united.
	   This two united power is integrated from gravity and antigravity.
	   Then this united power is zeta function.

	   $$ y =  1 + \tan^2 x = {1 \over {\cos x}} $$
	   $$ y^{'} = {{\cos^{'}x} \over {\cos^2 x}} $$
           $$ -{{\cos x} \over (\cos x)^{'}}(\sin x)^{'} = z_n $$
	   $$ {1 \over y^{'}} = z^2_n$$
           $$ z^n = - 2 e^{x \log x} $$

	   ガンマ関数の大域的積分多様体は、結局は、フェルマーの定理
	   だった。

         代数的計算手法のために$ \oplus L $を使っている。
そのために、冪乗計算と商代数の計算が、乗算で楽に見えるようになっている。
            微分幾何の量子化は、代数幾何の量子化の計算になっている。
  加減乗除が初等幾何であり、大域的微分と積分が、現代数学の代数計算の
  簡易での楽になる計算になっている。
  初等代数の計算は、
  $$ \bigoplus({i \hbar^{\nabla}})^{\oplus L} $$
   $$m \bigoplus({i \hbar^{\nabla}})^{\oplus L} +  
   n\bigoplus({i \hbar^{\nabla}})^{\oplus L} = 
   (m+n) \bigoplus({i \hbar^{\nabla}})^{\oplus L}   $$
   $$ m\bigoplus({i \hbar^{\nabla}})^{\oplus L} -  
   n\bigoplus({i \hbar^{\nabla}})^{\oplus L} = 
   (m-n) \bigoplus({i \hbar^{\nabla}})^{\oplus L}   $$
   $$ \bigoplus({i \hbar^{\nabla}})^{{\oplus L}^m} \times 
      \bigoplus({i \hbar^{\nabla}})^{{\oplus L}^n} =
   \bigoplus({i \hbar^{\nabla}})^{\oplus L^{m+n}} $$
  $$ {{\bigoplus({i \hbar^{\nabla}})^{\oplus L^m}}  \over  
  {\bigoplus({i \hbar^{\nabla}})^{\oplus L^n}}} =  
  \bigoplus({i \hbar^{\nabla}})^{\oplus L^{m\over n}} $$
  大域的計算での微分と積分は、
  $$ \left(\bigoplus({i \hbar^{\nabla}})^{\oplus L}\right)^{df}
  =  {\left(\bigoplus({i \hbar^{\nabla}})^{\oplus L}\right)^{
	  \bigoplus({i \hbar^{\nabla}})^{\oplus L}}}^{'}  $$
  $$ = \bigoplus({i \hbar^{\nabla}})^{\oplus L^{'}} $$
   $$ \int \bigoplus({i \hbar^{\nabla}})^{\oplus L} dx_m $$
   $$ \bigoplus({i \hbar^{\nabla}})^{\oplus L} = n $$
   大域的積分の代数幾何の量子化での計算は、ガンマ関数になっている。
   $$ {n^{L+1} \over {n+1}} = \int e^x x^{1 - t} dx $$
   代数幾何の量子化の因数分解は、加減乗除の式のm,nの
   組み合わせ多様体でのガンマ関数同士の計算になる。
   $$ \left(\bigoplus({i \hbar^{\nabla}})^{\oplus L}+m\right) 
   \left(\bigoplus({i \hbar^{\nabla}})^{\oplus L}+n\right) $$
  $$ = {n^{L+1} \over {n+1}} = t \int (1 - x)^n x^{m - 1} dx $$
  $$ \int x^{m-1}(1 - x)^{n-1} dx $$
  $$ \nabla ({i \hbar^{\nabla}})^{\oplus L}, \bigoplus ({i \hbar^{\nabla}})^{
	   \oplus L}, \square ({i \hbar^{\nabla}})^{\oplus L} $$
   $$ \boxtimes (i \hbar^{\nabla})|_{dx_m}^L, \boxplus (i \hbar^{\nabla}|_
   {dx_m}^L) $$



  $$ x^{{1 \over 2} + iy} = e^{x \log x} $$
  それゆえに、この定数はゼータ関数と微分幾何の量子化を因数にもつ
  素因子分解の式にもなっている。
  Therefore, this product also constructed with differential geometry
  of quantum level equation and zeta function.
  $$ = C $$ 
  そして、この関数はオイラーの定数から広中平祐定理による4重帰納法の
  オイラーの公式からの多様体積分へとのサーストン空間のスペクトラム関数
  ともなっている。
  And this function of Euler product respectrum of focus with
  Heisuke Hironaka manifold in four assembled of integral Euler equation.. 

  $$ C =  b_0 + \cfrac{c_1}{b_1 +
	  \cfrac{c_2}{b_2 +
	  \cfrac{c_3}{b_3 +
	  \cfrac{c_4}{b_4 + \cdots}}}} $$
  この方程式は指数による連分数としての役割も担っている。  
  This equation demanded with continued number of step function.
  $$ x^{{1 \over 2} + iy} = e^{x \log x} $$
  $$ (2.71828)^{2.828} = a^{a + b^{b + c^{c + d^{d \cdots }}}} $$
  $$ = \int e^f \cdot x^{1 - t}dx $$
  This represent is Gamma function in Euler product.
  Therefore this product is zeta function of global differential equation.

  $$ {d \over d f}F(v_{ij},h) = \int {e^{-f}\mathrm{dV}}
  [- \Delta v + \nabla_{i}\nabla_{j}
  v_{ij} - R_{ij}v_{ij} - v_{ij}\nabla_{i}\nabla_{j} + 2 <\nabla f,\nabla h>
  + (R + \nabla f^2)({v \over 2} - h)]$$
  $$ = {d \over d f}F = {2 {\int (R + \nabla_{i}\nabla_{j} f)^2} \over 
  { - (R + \Delta f)}}\mathrm{dm} $$
  これらの方程式は8種類の微分幾何の次元多様体として、そして、これらの
  多様体は曲平面による双対性をも生成している。 
  そして、このガウスの曲平面は、大域的微分多様体と微分幾何の量子化から
  素因数を形成してもいる。
  These equation are eight differential geometry of dimension calyement.
  And these calyement equation excluded into pair of dimension surface.
  This surface of Gauss function are global differential manifold,
  and differential geometry of quantum level.
  $$ F \ge {d \over d f}\int \int{1 \over (x \log x)^2}dx_m 
  + {d \over d f}\int \int{1 \over (y \log y)^{1 \over 2}}dy_m $$

 
  Differential geometry of quantum level constructed with Euler product
  and Gamma function being discatastrophed from continued fraction style.
  Gravity equation lend with varint of monotonicity of level expresented
  from gravity of letter varient formula.
  これらの方程式は基本群と大域的微分多様体をエスコートしていもいて、
  ヴェイユ予想がこのダランベルシアンの切断方程式たちから
  輸送のポートにもなっている。ベータ関数とガンマ関数がこれらのフォームラ
  の方程式を放出してもいて、結果、これらの方程式は広中平祐定理の複素
  多様体とグリーシャ教授によるペレルマン多様体からサポートされてもいる。
  この2名の教授は、一つは抽象理論をもう一方は具象理論を説明としている。
  These equation escourted into Global differential manifold and
  fundemental equation.
  Weil's theorem is imported from this equation in gravity of letters.
  Beta function and Gamma function are excluded with these formulas.
  These equation comontius from Heisuke Hironaga of complex manifold
  and Gresha professor of Perelman manifold.
  These two professos are one of abstract theorem and the other of
  visual manifold theorem.


     原子を分解したら重力になり、宇宙を分解したら原子になる。
     $$ \left(R^{\mu\nu} + {1 \over 2}\Lambda g_{ij}\right)^{\mu\nu} $$
     $$ = \kappa \left( T^{\mu\nu}_e + T^{\mu\nu}_y + T^{\mu\nu}_g +
     T^{\mu\nu}_h + \cdots \right) $$
     $$ {d \over dM}\left({\kappa T^{\mu\nu}} \right) \to {\bigoplus 
     ({i {\hbar^{{\nabla}}})}^{\oplus L}}, H \Psi = i \hbar \psi $$
     重力から原子への分解(1)
     $$ \bigoplus {({i \hbar^{\nabla})}^{\oplus L}} \to
          {d \over dM}\left(\kappa T^{\mu\nu} \right) $$ 
	  原子から重力への分解(2)
	  この(1)と(2)はベクトルの向きが違い、両方共、加群分解である。
    $$ ||ds^2|| = \lim_{x \to \infty} [\delta(x) 
    \int \int \int \pi \left( \sum_{k=0}^{\infty}
		{{}^n\sqrt{p},x \over n} \right)^{1 \over 2}d \tau]^{\mu\nu} $$
		種数を分解していったら、８種類の微分幾何、サーストン・
		ペレルマン多様体になり、心の影が宇宙を投影している
		代表が潮汐力、月の引力が地球に素粒子として影響を与えている。
    $$ ||ds^2|| = \int [D^2 \psi \otimes h \nu]d \tau $$
    として、D-braneが、この種数の分解の元になっている。
    $$ = {d \over d M}F $$
    と、微分基底に多様体Mをとっている。

    正規部分群による回転行列がベータ関数に行き着く。
    複素行列が虚数を因数にもつ素粒子方程式にもなる。

    8種類の微分幾何がそれぞれの固有の時間をもっている。
    M理論の種数を
    $$ {d \over d M}F $$
    で分解していったら、
    $$ = {d \over df}F = {{-2 \int {(R + \nabla_{i}\nabla_{j}) \over 
     {(\Delta + R)}}}}e^{-f}dV $$
     の各展開していった数式に分かれる。
     この数式がSeifert manifold 
     $$ = (x,y,z) / \Gamma $$
     である。

     重力場の積分記号の中で微分したら、積分記号の中は、
     電磁場方程式になり、大域的多様体となり、重力を分解したら、
     電磁場になり、その解が弱い力と強い力になる。
     QEDは、光をプリズムで分解したら、すべての幾何構造が生成される。
     そのプリズム自体が、サーストン・ペレルマン多様体である。
     $$ \int \left({\kappa T^{\mu\nu}}\right)^{\nabla}dx_m
     = \int e^{x \log x}\mathrm{div}(\mathrm{rot} E)dx_m $$
     $$ = e^{- x \log x} $$
     と、統一場理論の仕組みとなり、逆問題に行き着く。
     リサ・ランドール博士の誤差関数の共変微分の周りを覆っているのが、
     コイコライザーによる温度についての偏微分方程式とわかり、
     アーベル多様体が温度についての重力エネルギーを表しているのを、
     深谷賢治先生のローカルとグローバルな空間の味方から、
     アカシックレコードの子の代数幾何の量子化が、なぜ、量子レベルでの
     重力エネルギーが測れるか、熱エネルギーが重力エネルギーで表せるのを、
     Grisha Perelman博士とLisa Randall博士、深谷賢治先生、竹内薫先生たち
     から、わかりました。アカシックレコードはあるらしいです。

        オイラー数を多様体積分で施すと、
   ガンマ関数になる。この場合は、Frobeniusの定理によって、
   和と積の法則、確率論から導かれる結果、
   異次元と宇宙の関係に持っていける。
   オイラー数を多様体積分すると、ガンマ関数になる。
   オイラー数における頂点が、宇宙の特異点に対応して、
   辺がノルムに属して、面が一般相対性理論における無限の接線であり、
   $$ \text{F(頂点) - N(辺) + E(面)} = 2 $$
   $$ {(x,y,z) \over \Gamma} = \int{(- F\circ N + E)} + 
   (-F\circ E + N) + (N\circ E + F)dx_m $$
   $$ {{{\partial \over \partial f_S}(x,y,z)} \over {\Gamma}} $$
   ここで、オイラーの定数をも多様体積分にして、付け加えると、
   $$ = \int{({\int {1 \over x^s}dx} - \log x)}\mathrm{dvol} $$
   大域的多様体積分のガンマ関数に化ける。
   $$ = \int e^{-x}x^{1 - t}dx_m $$

   $AdS_5$多様体の方程式を、オイラー数の各部分単体に合わせると、
   $$ ||ds^2|| = e^{-2\pi T||\psi||}[\eta + \bar{h}(x)]dx^{\mu}dx^{\nu} 
   + T^2d^2 \psi $$
   開集合として、表せる。
   $$ = \mathcal{O}(x) $$
   この上の式たちから、ダランベルシアンは、$\square$と、オイラー数から演算子
   が決まっていて、共変微分は、$\nabla$で、双対被覆での単体での演算子で、
   決まっている。
   $$ \square = {{8 \pi G} \over c^4} T^{\mu\nu} $$
   $$ \nabla \psi = 4 \pi G \rho $$

   オイラー数を共変微分すると、偏微分方程式と同じ機能を、
   大域的積分多様体のオイラー数が、受け持っていて、
   このオイラー数の大域的多様体の共変微分で、ガンマ関数の大域的多様体
   が、
   ここで、オイラーの定数をも多様体積分にして、付け加えると、
   $$ = \int{({\int {1 \over x^s}dx} - \log x)}\mathrm{dvol} $$
   大域的多様体積分のガンマ関数に化ける。
   $$ = \int e^{-x}x^{1 - t}dx_m $$
   $$ = {d \over d f}F = \int \Gamma(\gamma)^{'}dx_m $$
   $$ = \int C dx_m = \int \Gamma dx_m \circ {d \over d \gamma}\Gamma $$
   $$ \le \int \Gamma dx_m + {d \over d \gamma}\Gamma $$
   $$ = e^{- \theta} + e^{i \theta} $$
   $$ = \square = 2(\cos (i x \log x) - i \sin (i x \log x))$$
   $$ = e^{-f} - e^{f} \le e^{-f} + e^{f} $$
   統一場理論は、彩さんのJones多項式に行き着く。

   Zata function escourt into Forgotfull of underlying modify on Sum
   of summative group, thiis Forgotfull summative group instented with
   Space ideal theorem from quantum level of deprivate space in aspect
   of quantum level of differential geometry,
   This theorem construct with Higgs field and Morse theory, moreover
   Hortshorn conjecture, fourth step deduction of theorem from 
   Euler product of integral manifold.
   This manfold revolte with global differential manifold in global topology.
   From these theorem inspirate from Forgotfull theorem in Space ideality
   theory, and this theory conclude into entropy of non exchange equation.
   Forgotten theory estimate with zeta function oneselves.
   Forgotten theory income with non relativity theory and this theory
   face to face without partner and partner of non relativity from
   ideal of space in revolution theory.
   All exist and non exist partner with non relativity escourt into 
   Forgotten theory, and this theory include with same underly group
   of all of partner with same distance ideal.

   $$ X \dotsi \to Y, Y \dotsi \to X $$
   $$ \bigoplus(i \hbar^{\nabla})^{\oplus L} = e^{-x \log x} $$
   This equation of quantum level of differential geomerty
   instimate with entorpy of non exchange equation.
   This equation means with low level of botton entropy in Space ideal of
   Forgotten theory. Moveover, this equation enterstein with all of exist
   theory from general relativity.
   No time and No Space of relativity, and monotonicity of magnetic component
   with gravity and antigravity, and this theory comment with partner of
   magnetic component theory included with being resulted from Forgotten
   theory.
   $$ \square \to \nabla |: x \to y, f(x) \to g(y), f^{-1}(x)xf(x) 
   \dotsi \to g(y) $$
   $$ \square \dotsi \to \nabla |: f(x) \cong g(y) $$
   $$ {d \over dt}g_{ij}(t) = - 2 R_{ij}, F^m_t = -2\int{{(R + \nabla_i \nabla
   _j f)} \over {(\Delta + f)}}dm $$
   $$ {d \over df}F = m(x) $$
   $$ F|_{\_} = {d \over df}F, F^m_t = e^{-f}dV $$
   $$ \bigoplus(i \hbar^{\nabla})^{\oplus L} = e^{-x \log x} $$
   These equation represent with Forgotten theory.
   空間概念の量子化は、代数幾何の量子化であり、忘却関数でもあり、
   忘却関手のカテゴリー論でもあり、簡潔に言うと、すべての相対性が
   機能すると、全て同型というコホモロジー論でのコイコライザーである。
   このコイコライザーが、空間概念の量子化である。
   すべては、ゼータ関数へと行き着く。
   ベータ関数は、誤差関数であり、宇宙の雑音として存在していると、
   異次元へと移行する。
   $$ \beta(p,q) \cong {\Gamma(p)\Gamma(q) \le \Gamma(p+q)}, 
   {{\Gamma(p)\Gamma(q)} \over \Gamma(p+q)} \le 1 $$
   $$ \int e^{-x}x^{1-t}dx_m = \int e^{-x}x^{1-t}dx\cdot e^{-\theta} $$
   $$ = \Gamma^{{\gamma}^{'}} $$
   $$ = e^{-x \log x} $$
   閉３次元多様体と３次元球面が、空間概念の量子化としての結果であり、
   この多様体が、無としてのベータ関数を形作っている３次元多様体である。
   宇宙と異次元合わせてが、空間概念の量子化での単体である。
   コイコライザーとは、共変微分によって対になっている変数に作用する写像であり、
   このコイコライザーで処理された変数同士を同型にさせる群が、
   同型のコイコライザーである。そして、最小の単体でもある。
   故に、宇宙と異次元合わせて、コイコライザーである。
   球対称である一般相対性理論の平坦な空間が、空間概念の量子化でもある。
   ビッグ・バンは、イコライザーであり、
   $$ \langle f,g\rangle \mapsto \langle d,d\rangle $$
   として、宇宙と異次元の曲平面を表している。
 
   宇宙には、最小値である273Kがあるから、リサ・ランドール博士は、
   平坦な宇宙である、球対称な一般相対性理論の解はないと言っているのが、
   わかりました。それで、アインシュタイン博士をおもって、ベクトル場で
   Warped Passagesにしているらしいです。
   光量子仮設までも書かれています。
   大学のテキストは、納得です。
   [Reference: 深谷賢治先生と加藤文元先生と竹内薫先生、S.マックレーン先生]

\end{document}

